\batchmode
\documentclass[12pt]{article}
\makeatletter

\usepackage{color,amsmath}

\pagestyle{empty}

\oddsidemargin0in
\textwidth6.5in
\topmargin0in
\textheight8.5in

\providecommand{\xin}{x_i}
\providecommand{\xia}{x_{i+1}}
\par\pagecolor[gray]{.7}


\makeatletter
\count@=\the\catcode`\_ \catcode`\_=8 
\newenvironment{tex2html_wrap}{}{} \catcode`\_=\count@
\makeatother
\let\mathon=$
\let\mathoff=$
\ifx\AtBeginDocument\undefined \newcommand{\AtBeginDocument}[1]{}\fi
\newbox\sizebox
\setlength{\hoffset}{0pt}\setlength{\voffset}{0pt}
\addtolength{\textheight}{\footskip}\setlength{\footskip}{0pt}
\addtolength{\textheight}{\topmargin}\setlength{\topmargin}{0pt}
\addtolength{\textheight}{\headheight}\setlength{\headheight}{0pt}
\addtolength{\textheight}{\headsep}\setlength{\headsep}{0pt}
\setlength{\textwidth}{349pt}
\newwrite\lthtmlwrite
\makeatletter
\let\realnormalsize=\normalsize
\global\topskip=2sp
\def\preveqno{}\let\real@float=\@float \let\realend@float=\end@float
\def\@float{\let\@savefreelist\@freelist\real@float}
\def\end@float{\realend@float\global\let\@freelist\@savefreelist}
\let\real@dbflt=\@dbflt \let\end@dblfloat=\end@float
\let\@largefloatcheck=\relax
\def\@dbflt{\let\@savefreelist\@freelist\real@dbflt}
\def\adjustnormalsize{\def\normalsize{\mathsurround=0pt \realnormalsize
 \parindent=0pt\abovedisplayskip=0pt\belowdisplayskip=0pt}\normalsize}%
\def\lthtmltypeout#1{{\let\protect\string\immediate\write\lthtmlwrite{#1}}}%
\newcommand\lthtmlhboxmathA{\adjustnormalsize\setbox\sizebox=\hbox\bgroup}%
\newcommand\lthtmlvboxmathA{\adjustnormalsize\setbox\sizebox=\vbox\bgroup%
 \let\ifinner=\iffalse }%
\newcommand\lthtmlboxmathZ{\@next\next\@currlist{}{\def\next{\voidb@x}}%
 \expandafter\box\next\egroup}%
\newcommand\lthtmlmathtype[1]{\def\lthtmlmathenv{#1}}%
\newcommand\lthtmllogmath{\lthtmltypeout{l2hSize %
:\lthtmlmathenv:\the\ht\sizebox::\the\dp\sizebox::\the\wd\sizebox.\preveqno}}%
\newcommand\lthtmlfigureA[1]{\let\@savefreelist\@freelist
       \lthtmlmathtype{#1}\lthtmlvboxmathA}%
\newcommand\lthtmlfigureZ{\lthtmlboxmathZ\lthtmllogmath\copy\sizebox
       \global\let\@freelist\@savefreelist}%
\newcommand\lthtmldisplayA[1]{\lthtmlmathtype{#1}\lthtmlvboxmathA}%
\newcommand\lthtmldisplayB[1]{\edef\preveqno{(\theequation)}%
  \lthtmldisplayA{#1}\let\@eqnnum\relax}%
\newcommand\lthtmldisplayZ{\lthtmlboxmathZ\lthtmllogmath\lthtmlsetmath}%
\newcommand\lthtmlinlinemathA[1]{\lthtmlmathtype{#1}\lthtmlhboxmathA  \vrule height1.5ex width0pt }%
\newcommand\lthtmlinlineA[1]{\lthtmlmathtype{#1}\lthtmlhboxmathA}%
\newcommand\lthtmlinlineZ{\egroup\expandafter\ifdim\dp\sizebox>0pt %
  \expandafter\centerinlinemath\fi\lthtmllogmath\lthtmlsetinline}
\newcommand\lthtmlinlinemathZ{\egroup\expandafter\ifdim\dp\sizebox>0pt %
  \expandafter\centerinlinemath\fi\lthtmllogmath\lthtmlsetmath}
\def\lthtmlsetinline{\hbox{\vrule width.1em\vtop{\vbox{%
  \kern.1em\copy\sizebox}\ifdim\dp\sizebox>0pt\kern.1em\else\kern.3pt\fi
  \ifdim\hsize>\wd\sizebox \hrule depth1pt\fi}}}
\def\lthtmlsetmath{\hbox{\vrule width.1em\vtop{\vbox{%
  \kern.1em\kern0.8 pt\hbox{\hglue.17em\copy\sizebox\hglue0.8 pt}}\kern.3pt%
  \ifdim\dp\sizebox>0pt\kern.1em\fi \kern0.8 pt%
  \ifdim\hsize>\wd\sizebox \hrule depth1pt\fi}}}
\def\centerinlinemath{%\dimen1=\ht\sizebox
  \dimen1=\ifdim\ht\sizebox<\dp\sizebox \dp\sizebox\else\ht\sizebox\fi
  \advance\dimen1by.5pt \vrule width0pt height\dimen1 depth\dimen1 
 \dp\sizebox=\dimen1\ht\sizebox=\dimen1\relax}

\def\lthtmlcheckvsize{\ifdim\ht\sizebox<\vsize\expandafter\vfill
  \else\expandafter\vss\fi}%
\makeatletter \tracingstats = 1 


\begin{document}
\pagestyle{empty}\thispagestyle{empty}%
\lthtmltypeout{latex2htmlLength hsize=\the\hsize}%
\lthtmltypeout{latex2htmlLength vsize=\the\vsize}%
\lthtmltypeout{latex2htmlLength hoffset=\the\hoffset}%
\lthtmltypeout{latex2htmlLength voffset=\the\voffset}%
\lthtmltypeout{latex2htmlLength topmargin=\the\topmargin}%
\lthtmltypeout{latex2htmlLength topskip=\the\topskip}%
\lthtmltypeout{latex2htmlLength headheight=\the\headheight}%
\lthtmltypeout{latex2htmlLength headsep=\the\headsep}%
\lthtmltypeout{latex2htmlLength parskip=\the\parskip}%
\lthtmltypeout{latex2htmlLength oddsidemargin=\the\oddsidemargin}%
\makeatletter
\if@twoside\lthtmltypeout{latex2htmlLength evensidemargin=\the\evensidemargin}%
\else\lthtmltypeout{latex2htmlLength evensidemargin=\the\oddsidemargin}\fi%
\makeatother
{\newpage\clearpage
\lthtmldisplayA{displaymath457}%
\begin{displaymath}|I - I_n| \leq \frac{C}{n^p}
\end{displaymath}%
\lthtmldisplayZ
\hfill\lthtmlcheckvsize\clearpage}

{\newpage\clearpage
\lthtmldisplayA{displaymath458}%
\begin{displaymath}|I_n - I_{n/2}| \leq \frac{D}{n^p}.
\end{displaymath}%
\lthtmldisplayZ
\hfill\lthtmlcheckvsize\clearpage}

{\newpage\clearpage
\lthtmlfigureA{alignstar18}%
\begin{align*}|I_n - I_{n/2}| & = |I_n \textcolor{blue}{- I + I} - I_{n/2}| &
\text{\textcolor{blue}{Add and subtract the same thing.}} \\
& \leq |I_n - I| \textcolor{blue}{+} | I - I_{n/2}| &
\text{\textcolor{blue}{Triangle inequality.}} \\
& \leq \textcolor{blue}{\frac{C}{n^p} + \frac{C 2^p}{n^p}}&
\text{\textcolor{blue}{Use the given bound.}} \\
& \leq \frac{C}{n^p}\textcolor{blue}{(1 + 2^p)}&
\text{\textcolor{blue}{So we see that $D = C(1+2^p)$ .}} \\
\end{align*}%
\lthtmlfigureZ
\hfill\lthtmlcheckvsize\clearpage}

{\newpage\clearpage
\lthtmldisplayA{displaymath459}%
\begin{displaymath}\int_{\frac{1}{10}}^1 \frac{\cos\left(\frac{\ln(x)}{x}\right)}{x} dx
\end{displaymath}%
\lthtmldisplayZ
\hfill\lthtmlcheckvsize\clearpage}

{\newpage\clearpage
\lthtmlinlinemathA{tex2html_wrap_inline513}%
$\epsilon$%
\lthtmlinlinemathZ
\hfill\lthtmlcheckvsize\clearpage}

{\newpage\clearpage
\lthtmldisplayA{displaymath460}%
\begin{displaymath}I \approx I_n = \frac{1}{2} f(x_1) h + \sum_{i=2}^{n} f(x_i) h +
\frac{1}{2} f(x_{n+1}) h 
\end{displaymath}%
\lthtmldisplayZ
\hfill\lthtmlcheckvsize\clearpage}

{\newpage\clearpage
\lthtmldisplayA{displaymath461}%
\begin{displaymath}|I - I_n| \leq \frac{M(b-a)^3}{12 n^2}
\end{displaymath}%
\lthtmldisplayZ
\hfill\lthtmlcheckvsize\clearpage}

{\newpage\clearpage
\lthtmlinlinemathA{tex2html_wrap_inline527}%
$|f''(x)| \leq M$%
\lthtmlinlinemathZ
\hfill\lthtmlcheckvsize\clearpage}

{\newpage\clearpage
\lthtmldisplayA{displaymath462}%
\begin{displaymath}\int_{x_i}^{x_{i+1}} f(x) dx = 
\int_{x_i}^{x_{i+1}} \left[ \frac{1}{2}f(x)+ \frac{1}{2}f(x) \right] dx 
\end{displaymath}%
\lthtmldisplayZ
\hfill\lthtmlcheckvsize\clearpage}

{\newpage\clearpage
\lthtmlfigureA{alignstar75}%
\begin{align*}\frac{1}{2}f(x) & =
\frac{1}{2} \left[
f(x_i) + f'(x_i)(x-x_i) + \frac{1}{2} f''(x_i)(x-x_i)^2 +
\frac{1}{3!}  f'''(\xi_1)(x-x_i)^3
\right] \\
\frac{1}{2}f(x) & =
\frac{1}{2} \left[
f(x_{i+1}) + f'(x_{i+1})(x-x_{i+1}) + \frac{1}{2} f''(x_{i+1})(x-x_{i+1})^2 +
\frac{1}{3!}  f'''(\xi_1)(x-x_{i+1})^3
\right] \\
\end{align*}%
\lthtmlfigureZ
\hfill\lthtmlcheckvsize\clearpage}

{\newpage\clearpage
\lthtmlfigureA{alignstar93}%
\begin{align*}\int_{x_i}^{x_{i+1}} (x-x_i) dx, & = \frac{1}{2} h^2, & 
\int_{x_i}^{x_{i+1}} (x-x_{i+1}) dx, & = -\frac{1}{2} h^2, \\
\int_{x_i}^{x_{i+1}} (x-x_i)^2 dx, & = \frac{1}{3} h^3, &
\int_{x_i}^{x_{i+1}} (x-x_{i+1})^2 dx, & = \frac{1}{3} h^3, \\
\int_{x_i}^{x_{i+1}} (x-x_i)^3 dx, & = \frac{1}{4} h^4, &
\int_{x_i}^{x_{i+1}} (x-x_{i+1})^3 dx, & = -\frac{1}{4} h^4.
\end{align*}%
\lthtmlfigureZ
\hfill\lthtmlcheckvsize\clearpage}

{\newpage\clearpage
\lthtmlfigureA{multline119}%
\begin{multline}\int_{x_i}^{x_{i+1}} \frac{1}{2}f(x)+ \frac{1}{2}f(x) dx  = 
\frac{1}{2} h \left[f(x_i) + f(x_{i+1}) \right] + \\
\frac{1}{4} h^2 
[f'(x_i) - f'(x_{i+1})] + 
\frac{1}{12} h^3 [f''(x_i) + f''(x_{i+1})] + \\
\frac{1}{12} \int_{x_i}^{x_{i+1}} 
[f'''(\xi_1)(x-x_i)^3 + 
f'''(\xi_2)(x-x_{i+1})^3]dx. 
\end{multline}%
\lthtmlfigureZ
\hfill\lthtmlcheckvsize\clearpage}

{\newpage\clearpage
\lthtmlfigureA{align137}%
\begin{align}f'(x_{i+1}) & = f'(x_i) + h f''(x_i) + \frac{1}{2} h^2 f'''(\xi_3), \\
f''(x_{i+1}) & = f''(x_i) + h f'''(\xi_4).
\end{align}%
\lthtmlfigureZ
\hfill\lthtmlcheckvsize\clearpage}

{\newpage\clearpage
\lthtmlfigureA{multlinestar141}%
\begin{multline*}\int_{x_i}^{x_{i+1}} \frac{1}{2}f(x)+ \frac{1}{2}f(x) dx  = 
\frac{1}{2} h \left[f(x_i) + f(x_{i+1}) \right] + \\
-\frac{1}{4} h^3 f''(x_i) - \frac{1}{8} h^4 f'''(\xi_3)  +
\frac{1}{6} h^3 f''(x_i) + \frac{1}{6} h^4 f''(\xi_4)\\
\frac{1}{12} \int_{x_i}^{x_{i+1}} 
[f'''(\xi_1)(x-x_i)^3 + 
f'''(\xi_2)(x-x_{i+1})^3]dx. 
\end{multline*}%
\lthtmlfigureZ
\hfill\lthtmlcheckvsize\clearpage}

{\newpage\clearpage
\lthtmldisplayA{displaymath463}%
\begin{displaymath}\int_{x_i}^{x_{i+1}} f(x) dx  = \frac{1}{2} h \left[f(x_i) + f(x_{i+1})
\right] +
-\frac{1}{12} h^3 f''(x_i) + O(h^4).
\end{displaymath}%
\lthtmldisplayZ
\hfill\lthtmlcheckvsize\clearpage}

{\newpage\clearpage
\lthtmldisplayA{displaymath464}%
\begin{displaymath}\left|\int_{x_i}^{x_{i+1}} f(x) dx  - 
\frac{1}{2} h \left[f(x_i) + f(x_{i+1})
\right]\right|  \leq
\frac{1}{12} h^3 |f''(x_i)| + O(h^4),
\end{displaymath}%
\lthtmldisplayZ
\hfill\lthtmlcheckvsize\clearpage}

{\newpage\clearpage
\lthtmldisplayA{displaymath465}%
\begin{displaymath}p(x) = c_0 + c_1 x + \frac{c_2}{2} x^2.
\end{displaymath}%
\lthtmldisplayZ
\hfill\lthtmlcheckvsize\clearpage}

{\newpage\clearpage
\lthtmldisplayA{displaymath466}%
\begin{displaymath}I = \int_0^L p(x) dx = c_0 L + c_1 \frac{L^2}{2} + c_2 \frac{L^3}{6},
\end{displaymath}%
\lthtmldisplayZ
\hfill\lthtmlcheckvsize\clearpage}

{\newpage\clearpage
\lthtmldisplayA{displaymath467}%
\begin{displaymath}T_N = \frac{h}{2} \left[p(0) + p(L)\right] + h \sum_{i=2}^{N} p(x_i).
\end{displaymath}%
\lthtmldisplayZ
\hfill\lthtmlcheckvsize\clearpage}

{\newpage\clearpage
\lthtmlinlinemathA{tex2html_wrap_inline539}%
$h = \frac{L}{N}$%
\lthtmlinlinemathZ
\hfill\lthtmlcheckvsize\clearpage}

{\newpage\clearpage
\lthtmldisplayA{displaymath468}%
\begin{displaymath}T_N = \frac{L}{2N} \left[p(0) + p(L)\right] + 
\frac{L}{N} \sum_{i=1}^{N-1} p\left(\frac{iL}{N}\right).
\end{displaymath}%
\lthtmldisplayZ
\hfill\lthtmlcheckvsize\clearpage}

{\newpage\clearpage
\lthtmlfigureA{alignstar216}%
\begin{align*}T_N & = \frac{L}{2N} 
\left[c_0 + c_0 + c_1 L + \frac{1}{2} c_2 L^2
\right] +
\frac{L}{N} \sum_{i=1}^{N-1}
\left[
c_0 + c_1 \frac{iL}{N} + \frac{c_2}{2} \left(\frac{iL}{N}\right)^2
\right] \\
& = c_0 \left[\frac{L}{N} + \frac{L (N-1)}{N}
\right] +
c_1 \left[
\frac{L^2}{2N} + \frac{L^2}{N^2} \sum_{i=1}^{N-1} i
\right] +
c_2 \left[ \frac{L^3}{4 N} + \frac{L^3}{2 N^3} \sum_{i=1}^{N-1} i^2
\right] \\
& =
c_0 L + 
c_1 \left[
\frac{L^2}{2N} + \frac{L^2}{N^2} \frac{N(N-1)}{2}
\right] +
c_2 \left[ \frac{L^3}{4 N} + \frac{L^3}{2 N^3} \frac{N (N-1) (2N-1)}{6}
\right] \\
& =
c_0 L + 
c_1 \frac{L^2}{2} +
c_2 \left[ \frac{L^3}{6} + \frac{L^3}{12 N^2} 
\right]
\end{align*}%
\lthtmlfigureZ
\hfill\lthtmlcheckvsize\clearpage}

{\newpage\clearpage
\lthtmldisplayA{displaymath469}%
\begin{displaymath}I - T_N = -\frac{c_2 L^3}{12 N^2} 
\end{displaymath}%
\lthtmldisplayZ
\hfill\lthtmlcheckvsize\clearpage}

{\newpage\clearpage
\lthtmlinlinemathA{tex2html_wrap_inline65}%
$\mbox{$\bullet$}$%
\lthtmlinlinemathZ
\hfill\lthtmlcheckvsize\clearpage}

{\newpage\clearpage
\lthtmlinlinemathA{tex2html_wrap_inline555}%
$p'(\xi)$%
\lthtmlinlinemathZ
\hfill\lthtmlcheckvsize\clearpage}

{\newpage\clearpage
\lthtmlinlinemathA{tex2html_wrap_inline557}%
$\xi$%
\lthtmlinlinemathZ
\hfill\lthtmlcheckvsize\clearpage}

{\newpage\clearpage
\lthtmldisplayA{displaymath470}%
\begin{displaymath}p(0) = p_0, \ \ \ \ p(1) = p_1, \ \ \ \ p'(\xi) = p_\xi'.
\end{displaymath}%
\lthtmldisplayZ
\hfill\lthtmlcheckvsize\clearpage}

{\newpage\clearpage
\lthtmlinlinemathA{tex2html_wrap_inline559}%
$p_\xi'$%
\lthtmlinlinemathZ
\hfill\lthtmlcheckvsize\clearpage}

{\newpage\clearpage
\lthtmlfigureA{alignstar272}%
\begin{align*}c_0 & = p_0, \\
c_0 + c_1 + c_2 & = p_1, \\
c_1 + 2 c_2 \xi & = p_\xi'.
\end{align*}%
\lthtmlfigureZ
\hfill\lthtmlcheckvsize\clearpage}

{\newpage\clearpage
\lthtmldisplayA{displaymath472}%
\begin{displaymath}\left[
\begin{array}{ccc}
1 & 0 & 0 \\
1 & 1 & 1 \\
0 & 1 & 2 \xi
\end{array}
\right]
\left[
\begin{array}{c}
c_0 \\
c_1 \\
c_2
\end{array}
\right] =
\left[
\begin{array}{c}
p_0 \\
p_1 \\
p_\xi'
\end{array}
\right]
\end{displaymath}%
\lthtmldisplayZ
\hfill\lthtmlcheckvsize\clearpage}

{\newpage\clearpage
\lthtmlinlinemathA{tex2html_wrap_inline561}%
$2\xi - 1$%
\lthtmlinlinemathZ
\hfill\lthtmlcheckvsize\clearpage}

{\newpage\clearpage
\lthtmlinlinemathA{tex2html_wrap_inline563}%
$\xi = \frac{1}{2}$%
\lthtmlinlinemathZ
\hfill\lthtmlcheckvsize\clearpage}

{\newpage\clearpage
\lthtmlinlinemathA{tex2html_wrap_inline565}%
$\xi
\neq \frac{1}{2}$%
\lthtmlinlinemathZ
\hfill\lthtmlcheckvsize\clearpage}

{\newpage\clearpage
\lthtmlfigureA{alignstar289}%
\begin{align*}c_0 & = p_0, \\
c_0 + c_1 + c_2 + c_3& = p_1, \\
c_1 + c_2 + \frac{3}{4} c_3 & = p_\xi'.
\end{align*}%
\lthtmlfigureZ
\hfill\lthtmlcheckvsize\clearpage}

{\newpage\clearpage
\lthtmldisplayA{displaymath474}%
\begin{displaymath}\left[
\begin{array}{ccc}
1 & 0 & 0 \\
1 & 1 & 0 \\
0 & 1 & \frac{3}{4} \xi
\end{array}
\right]
\left[
\begin{array}{c}
c_0 \\
d \\
c_3
\end{array}
\right] =
\left[
\begin{array}{c}
p_0 \\
p_1 \\
p_\xi'
\end{array}
\right].
\end{displaymath}%
\lthtmldisplayZ
\hfill\lthtmlcheckvsize\clearpage}

{\newpage\clearpage
\lthtmlinlinemathA{tex2html_wrap_inline587}%
$\frac{3}{4}$%
\lthtmlinlinemathZ
\hfill\lthtmlcheckvsize\clearpage}

{\newpage\clearpage
\lthtmldisplayA{displaymath475}%
\begin{displaymath}\int_0^1 f(x) dx
\end{displaymath}%
\lthtmldisplayZ
\hfill\lthtmlcheckvsize\clearpage}

{\newpage\clearpage
\lthtmlinlinemathA{tex2html_wrap_inline601}%
$0, \frac{1}{3}, \frac{2}{3}, 1$%
\lthtmlinlinemathZ
\hfill\lthtmlcheckvsize\clearpage}

{\newpage\clearpage
\lthtmlinlinemathA{tex2html_wrap_inline607}%
$3 \left[f\left(\frac{1}{3}\right)-f(0)\right]$%
\lthtmlinlinemathZ
\hfill\lthtmlcheckvsize\clearpage}

{\newpage\clearpage
\lthtmlinlinemathA{tex2html_wrap_inline609}%
$\frac{1}{3}$%
\lthtmlinlinemathZ
\hfill\lthtmlcheckvsize\clearpage}

{\newpage\clearpage
\lthtmlinlinemathA{tex2html_wrap_inline611}%
$f\left(\frac{1}{3}\right)$%
\lthtmlinlinemathZ
\hfill\lthtmlcheckvsize\clearpage}

{\newpage\clearpage
\lthtmlinlinemathA{tex2html_wrap_inline613}%
$\frac{9}{2} \left[f\left(\frac{2}{3}\right)-2
f\left(\frac{1}{3}\right)+
f(0)\right]$%
\lthtmlinlinemathZ
\hfill\lthtmlcheckvsize\clearpage}

{\newpage\clearpage
\lthtmlinlinemathA{tex2html_wrap_inline615}%
$3 \left[f\left(\frac{2}{3}\right)-f\left(\frac{1}{3}\right)\right]$%
\lthtmlinlinemathZ
\hfill\lthtmlcheckvsize\clearpage}

{\newpage\clearpage
\lthtmlinlinemathA{tex2html_wrap_inline617}%
$\frac{9}{2}\left[f(1)-3 f\left(\frac{2}{3}\right)+
3 f\left(\frac{1}{3}\right)-f(0)\right]$%
\lthtmlinlinemathZ
\hfill\lthtmlcheckvsize\clearpage}

{\newpage\clearpage
\lthtmlinlinemathA{tex2html_wrap_inline619}%
$\frac{2}{3}$%
\lthtmlinlinemathZ
\hfill\lthtmlcheckvsize\clearpage}

{\newpage\clearpage
\lthtmlinlinemathA{tex2html_wrap_inline621}%
$f\left(\frac{2}{3}\right)$%
\lthtmlinlinemathZ
\hfill\lthtmlcheckvsize\clearpage}

{\newpage\clearpage
\lthtmlinlinemathA{tex2html_wrap_inline623}%
$\frac{9}{2}\left[f(1)-
2 f\left(\frac{2}{3}\right)+f\left(\frac{1}{3}\right)\right]$%
\lthtmlinlinemathZ
\hfill\lthtmlcheckvsize\clearpage}

{\newpage\clearpage
\lthtmlinlinemathA{tex2html_wrap_inline625}%
$3 \left[f(1)-f\left(\frac{2}{3}\right)\right]$%
\lthtmlinlinemathZ
\hfill\lthtmlcheckvsize\clearpage}

{\newpage\clearpage
\lthtmlfigureA{multlinestar350}%
\begin{multline*}p(x) =
f(0) + 3 \left[f\left(\frac{1}{3}\right)-f(0)\right] x +
\frac{9}{2} \left[f\left(\frac{2}{3}\right)-2
f\left(\frac{1}{3}\right)+
f(0)\right] x \left(x-\frac{1}{3}\right) +
\\
\frac{9}{2}\left[f(1)-3 f\left(\frac{2}{3}\right)+
3 f\left(\frac{1}{3}\right)-f(0)\right]  
x  \left(x-\frac{1}{3}\right) \left(x-\frac{2}{3}\right).
\end{multline*}%
\lthtmlfigureZ
\hfill\lthtmlcheckvsize\clearpage}

{\newpage\clearpage
\lthtmlfigureA{multlinestar372}%
\begin{multline*}\int_0^1 p(x) dx =
f(0) + \frac{3}{2} \left[f\left(\frac{1}{3}\right)-f(0)\right]  +
\frac{3}{4} \left[f\left(\frac{2}{3}\right)-2
f\left(\frac{1}{3}\right)+
f(0)\right]  +
\\
\frac{1}{8}\left[f(1)-3 f\left(\frac{2}{3}\right)+
3 f\left(\frac{1}{3}\right)-f(0)\right].
\end{multline*}%
\lthtmlfigureZ
\hfill\lthtmlcheckvsize\clearpage}

{\newpage\clearpage
\lthtmldisplayA{displaymath476}%
\begin{displaymath}\int_0^1 p(x) dx = \frac{1}{8} f(0) +
\frac{3}{8} f\left(\frac{1}{3}\right) +
\frac{3}{8} f\left(\frac{2}{3}\right) +
\frac{1}{8} f(1).
\end{displaymath}%
\lthtmldisplayZ
\hfill\lthtmlcheckvsize\clearpage}

{\newpage\clearpage
\lthtmlinlinemathA{tex2html_wrap_inline633}%
$\frac{3}{8}$%
\lthtmlinlinemathZ
\hfill\lthtmlcheckvsize\clearpage}


\end{document}
